% [原文 p.7]
\[
\begin{array}{rcl}
 \fM^{0,c}_A &=& \text{复化 Kähler 类 }[\omega-iB]\in H^2(X,\R/2\pi i\Z)\text{ 的空间},\\
 \fM^{0,c}_B &=& X \text{ 的复结构的模空间}.
\end{array}
\]
$\fM^{0,c}_{A/B}$ 是 $\fM_{A/B}$ 的一个子流形，对应于经典系统的边缘形变。当 $c_1(X)=0$ 时，
空间 $\fM^{0,c}_A\times\fM^{0,c}_B$ 可等同于红外 SCFT 模空间的一个开子集。当 $c_1(X)\neq 0$ 时，
在 $\fM^{0,c}_A$ 上存在沿 $c_1(X)$ 方向的 RG 流，
而沿 $ic_1(X)$ 方向的平移则被轴转动吸收。
% [Tutor Note: 印刷 p.7 原文此处作 "there is an RG low"，其中 "low" 应为 "flow"；同句 "rotaion" 应为 "rotation"。译文按正确含义处理，未改动原文。]

由于该模型具有带荷整性的矢量 $U(1)$ R-对称性，A-扭转总是可行的；相应的拓扑弦论在数学中称为
\emph{Gromov--Witten 理论}。B-扭转可行当且仅当 $X$ 是 Calabi--Yau 流形，即典范丛 $K_X$
作为全纯线丛是平凡的。

\para{Landau--Ginzburg 模型}

设 $W(x)$ 是 $N$ 个变量 $x=(x^1,\dots,x^N)$ 的复系数多项式，且只有孤立临界点。则存在一个以
$W(x)$ 为超势的二维 $(2,2)$ 超对称 QFT，称为 \emph{Landau--Ginzburg 模型}。

其场变量由复标量场 $x^1,\dots,x^N$ 与 Dirac 费米子 $\psi^1,\dots,\psi^N$ 组成。在 Minkowski
时空上，拉氏量为
\begin{align}
  \cL &= \sum_{i=1}^N\Big(|\partial_0x^i|^2-|\partial_1x^i|^2
        +i\bar\psi^i_-\big(\partial_0+\partial_1\big)\psi^i_-
        +i\bar\psi^i_+\big(\partial_0-\partial_1\big)\psi^i_+\Big)\notag\\
      &\quad -\frac14\sum_{i=1}^N|\partial_iW(x)|^2
        -\frac12\sum_{i,j=1}^N\Big(\psi^i_+\psi^j_-\,\partial_i\partial_jW(x)+\text{c.c.}\Big),
        \tag{2.17}
\end{align}
而超对称变换为
\begin{equation}
  \delta x^i=\epsilon_+\psi^i_--\epsilon_-\psi^i_+, \tag{2.18}
\end{equation}
\begin{equation}
  \delta\psi^i_\pm=\pm i\bar\epsilon_\mp(\partial_0\pm\partial_1)x^i
   -\frac12\epsilon_\pm\overline{\partial_iW(x)}. \tag{2.19}
\end{equation}

该模型总具有带荷整性的轴 $U(1)$ R-对称性。矢量 R-对称性存在当且仅当 $W(x)$ 是拟齐次的，即
必要时作坐标变换后，存在一组数 $R=(R_1,\dots,R_N)$ 使得 $W(\lambda^Rx)=\lambda^2W(x)$，其中
$\lambda^Rx=(\lambda^{R_1}x_1,\dots,\lambda^{R_N}x_N)$。此时理论在红外极限下流向一个
$\hat c=\tr(1-R)$ 的 SCFT。该模型的手征环等为
% [原文 p.8]
\[
\begin{array}{rcl@{\qquad}l}
 \cR_A &=& ?,\\
 \cR_B &=& \mathrm{Jac}(W) & \text{Jacobi 环}\ \C[x_1,\dots,x_N]/(\partial_1W,\dots,\partial_NW),\\
 \cC_A &=& \mathrm{Fuk}(W) & \text{由 }W\text{ 控制的 Fukaya 范畴},\\
 \cC_B &=& \mathrm{MF}(W) & W\text{ 的矩阵分解范畴}\\
 \fM_A &=& ?,\\
 \fM_B &=& \multicolumn{2}{l}{\text{$W$ 的泛形变的模空间},}
\end{array}
\]
讲演者目前还不知道 $\cR_A$ 与 $\fM_A$ 究竟是什么。

由于该模型具有带荷整性的轴 $U(1)$ R-对称性，B-扭转总是可行的。相应的拓扑弦论在树图（亏格零）
层面与 K.~Saito 的原始形式理论密切相关。当 $W$ 是 Morse 函数时，A.~Givental 提出了构造高亏格
振幅的方案，而 C.~Teleman 证明了它们满足拓扑弦论的数学公理。

当 $W$ 是拟齐次时，还存在矢量 R-对称性。然而它不具荷整性，因而 A-扭转不可行。这个问题可以通过
轨形化来补救：用 $W(x)$ 的一个有限对称群 $\Gamma\subset GL(N,\C)$ 来规范该理论，并须使
$\ee^{\pi iR}$ 成为 $\Gamma$ 的元素。于是对规范不变场荷整性成立，A-扭转变得可行。相应的拓扑弦论
称为 \emph{FJRW 理论}。

\subsection{镜像对称}
\label{sec:mirror}

二维 $(2,2)$ 超对称代数有一个自同构：$Q_-\leftrightarrow\overline Q_-$、$F_V\leftrightarrow F_A$，
其余生成元保持不变。若两个二维 $(2,2)$ 超对称 QFT 之间存在一个同构，使得在该同构下超对称生成元
经历上述自同构，则称它们互为\emph{镜像}。镜像对称有若干直接的推论：一个理论的环 $\cR_A$ 同构于
其镜像的环 $\cR_B$；一个理论的范畴 $\cC_A$ 等价于其镜像的范畴 $\cC_B$；一个理论的参数空间
$\fM_A$ 同构于其镜像的参数空间 $\fM_B$；一个理论的拓扑 A-模型（A-扭转后的模型或相应的拓扑弦论）
同构于另一个理论的拓扑 B-模型。

镜像对称最著名的例子是以 Calabi--Yau 流形（比如 $X$ 与 $Y$）为靶的 sigma 模型之间的镜像对称。
作为上述推论的一部分，我们有 Hodge 数之间的关系
$h^{p,q}(X)=h^{n-p,q}(Y)$，其中 $n=\dim X=\dim Y$。其他众所周知的例子包括：以非 Calabi--%
% [原文 p.9]
Yau 靶为靶的 sigma 模型与 Landau--Ginzburg 模型之间的镜像对称，以及 Landau--Ginzburg 轨形
之间的镜像对称。这些例子中的一些推论已在数学上得到证明。

% [Tutor Note: p.8 的 "R_A = ?" 与 "M_A = ?" 是作者刻意留下的问号：原文紧接着写 "The lecturer
%  does not know what R_A and M_A are at this moment."，说明截至 2016 年 A-扭转侧的 LG 手征环与
% 参数空间尚无已知描述。译文照排问号，不补内容。]
% [Tutor Note: p.8 原文作 "Jocobi ring"，应为 "Jacobi ring"（雅可比环）之笔误；C[x_1,...,x_N]/(∂_1W,...,∂_NW)
% 的正确描述是多项式环模去由 W 的所有一阶偏导生成的理想。译文按正确含义作 "Jacobi 环" 并在正文保留原式。]
% [Tutor Note: 关于中心荷：原文在 §2.3 引入 $\hat c=c/3$（而非 $c=3\hat c$ 的等价说法）；在 $(2,2)$ SCFT 中
% $\hat c$ 即 R-对称性的 't Hooft 异常系数。对 LG 模型 $\hat c=\tr(1-R)$ 可用 R-荷迹验算：超势 W 的 R-荷为 2，
% 故 $\partial_iW$ 的 R-荷为 $2-R_i$；超势项 $\psi^i_+\psi^j_-\partial_i\partial_jW$ 的总 R-荷为 2，恰好给出
% $R_i+R_j+(2-R_i-R_j)=2$，与 $W(\lambda^Rx)=\lambda^2W(x)$ 的拟齐次性自洽，可作交叉检验。]
